\section{Which operations prepare and read the pair?}
\label{sec:pair-currents}

Sections~\ref{sec:preparation-evolution} and~\ref{sec:preparation-effective-dynamics} specified the pair preparation by an orbital rotation and its detector by an occupation imbalance. Both are built from three one-body operators. This section identifies them, shows that $\Span\{g,e,m\}$ is their representation of spin one in the orbital index, and asks which drives and mixers leave the ground in place. The answer separates what the pair Hamiltonian supplies from what the apparatus has to add. In the trap of Section~\ref{subsec:orbital-apparatus} the same three operators are the angular momentum and the two quadrupoles.

\subsection{Three currents and the space they act on}
\label{subsec:pair-currents-triplet}

With the modes of Section~\ref{subsec:pair-grounds-excitations}, put $N_c=\sum_ic_i^*c_i$, $N_d=\sum_id_i^*d_i$ and
\begin{equation}
 T=\sum_{i=1}^2c_i^*d_i,\qquad X=\tfrac12(T+T^*),\qquad Y=\tfrac1{2i}(T-T^*),\qquad Z=\tfrac12(N_c-N_d).
 \label{eq:pair-currents}
\end{equation}
Each conserves the particle number and acts in the same way on the two internal labels $i=1,2$. The canonical anticommutation relations give $[X,Y]=iZ$ with its cyclic permutations, and
\begin{equation}
 T^2=2P_c^*P_d .
 \label{eq:pair-currents-square}
\end{equation}
Indeed the terms of $T^2$ with equal labels vanish because $c_i^*c_i^*=0$, and each of the two with different labels equals $c_1^*c_2^*d_2d_1$. The pair transfer is therefore two successive single transfers, $P_c^*P_d+P_d^*P_c=\tfrac12(T^2+T^{*2})=X^2-Y^2$. On the two-particle space,
\begin{equation}
 \begin{aligned}
 Xg&=m, & Xe&=0, & Xm&=g,\\
 Yg&=0, & Ye&=im, & Ym&=-ie,\\
 Zg&=e, & Ze&=g, & Zm&=0,
 \end{aligned}
 \label{eq:pair-currents-action}
\end{equation}
and all three currents annihilate $W$. Thus $\Span\{g,e,m\}$ carries the representation of spin one of the currents, with $X^2+Y^2+Z^2=2$ there, and $W$ consists of three invariant vectors. We call this spin \emph{orbital}, because the currents act on the orbital index $c,d$ and leave the internal label alone. Each of $g$, $e$ and $m$ is the vector annihilated by one current, namely $Y$, $X$ and $Z$. Since $P_c^*P_c+P_d^*P_d$ and $Z^2$ agree on $\mathcal F_2$, the symmetric family of Section~\ref{subsec:pair-grounds-excitations} is
\begin{equation}
 H_s=-sZ^2-j\,(X^2-Y^2)\quad\text{on }\mathcal F_2,\qquad
 [Y,H_s]=i\,(j-s)\,(XZ+ZX).
 \label{eq:pair-currents-hamiltonian}
\end{equation}
The commutator is nonzero for $s\ne j$, since $(XZ+ZX)e=m$. The coupling $j$ is the constant $J=4\kappa$ of Singh \cite[Eq.~(3.11)]{Singh2026PairTransfer}, computed there for a specified scalar trace and matrix mode. With another normalization of the mode the coefficient changes, to $4\kappa/N$ in the example given there, so a numerical value of $j$ refers to that convention.

The rotation $\kappa(\varphi)$ of Section~\ref{subsec:symmetry-pair} is generated by $Y$. From $[Y,c_i^*]=\tfrac i2d_i^*$, $[Y,d_i^*]=-\tfrac i2c_i^*$ and $Y\ket{\mathrm{vac}}=0$,
\begin{equation}
 \kappa(\varphi)=e^{-2i\varphi Y}.
 \label{eq:pair-currents-rotation}
\end{equation}
Hence $A_\theta=\kappa(\theta)A_0$. The rotation fixes $g$ because $Yg=0$, and by Eq.~\eqref{eq:pair-currents-action} it turns $e$ into $b=u_\theta e+v_\theta m$. The measured imbalance is the rotated current,
\begin{equation}
 Z_\theta=\kappa(\theta)Z\kappa(\theta)^*=u_\theta Z+v_\theta X,\qquad
 Z_\theta g=b,\quad Z_\theta b=g,\quad Z_\theta d=0 .
 \label{eq:pair-currents-readout}
\end{equation}
On $\Span\{g,e,m\}$ this gives $Z_\theta^2=A_\theta A_\theta^*$ and $I-Z_\theta^2=\Pi_d$. The zero outcome of the imbalance is the projection onto the line that the preparation misses, and for a state $\rho$ supported on this space
\begin{equation}
 \Tr(\rho\,\Pi_d)=\Pr\nolimits_\rho(Z_\theta=0)=1-\Tr(\rho Z_\theta^2).
 \label{eq:pair-currents-zero}
\end{equation}
This identification is confined to $\Span\{g,e,m\}$. On $\mathcal F_2$ the currents annihilate $W$, so the zero outcome has the projection $\Pi_d+\Pi_W$ of Eq.~\eqref{eq:pair-ambient-zero}. A state on $d$ and a state in $W$ both give zero with certainty. On $\Span\{g,e,m\}$ the frequency of the outcome zero, or equivalently the second moment of the imbalance, determines the weight on $d$. The first moment does not. It vanishes for the state $\Pi_d$ and also for the equal mixture of the two eigenstates with outcomes $\pm1$. Their probabilities of zero are one and zero.

\subsection{Which drives keep the ground in place?}
\label{subsec:pair-currents-drives}

A Hermitian one-body operator that conserves the particle number and acts identically on the two labels has a Hermitian $2\times2$ matrix in the orbital index. It is a real combination of $N_c+N_d$, $X$, $Y$ and $Z$, and $N_c+N_d=2$ on $\mathcal F_2$. Consider therefore the drive $\alpha X+\beta Y+\gamma Z$, with real coefficients that may depend continuously on time. By Eq.~\eqref{eq:pair-currents-action} it sends $g$ to $\alpha m+\gamma e$. If the evolving ground stays on its ray throughout the pulse, differentiating shows that the drive maps $g$ into $\C g$ at every time. Conversely a multiple of $Y$ annihilates $g$. Hence
\begin{equation}
 \text{the ground ray is invariant throughout the pulse}
 \quad\Longleftrightarrow\quad \alpha=\gamma=0 .
 \label{eq:pair-currents-drive}
\end{equation}
The drives that keep the ground are the multiples of $Y$, and by Eq.~\eqref{eq:pair-currents-rotation} they produce the orbital rotations. Access to this drive and a fixed angle are assumptions about the apparatus. Singh's canonical component supplies $H_0$ and its ground. It does not supply the drive.

The pair component cannot replace the drive. Let $\Pi_c=(-1)^{N_c}$. The Hamiltonian $H_s$, the two Hermitian quadratures of $P_c^*P_d$, the imbalance $Z$ and the self-channels $P_c^*P_c$, $P_d^*P_d$ commute with $\Pi_c$ and map the pair plane $\Span\{C,D_p\}$ into itself. On that plane $\Pi_c=1$, while $\Pi_cm=-m$. No sequence of these controls takes a vector of the pair plane to one with a component on $m$. The currents $X$ and $Y$ change $N_c$ by one and anticommute with $\Pi_c$. A one-body mixer is therefore an operation in addition to the pair component. This is a restriction on the stated family of controls. It does not exclude a larger system with further sources.

The input generator of Section~\ref{subsec:pair-grounds-excitations} is the restriction of $H_s$ to the original pair plane $\Span\{C,D_p\}$, which is $\operatorname{diag}(-(s+j),\,j-s)$ in the basis $g_\Psi,e_\Psi$. The compression to the rotated plane $\Span\{g,b\}$ is different,
\begin{equation}
 A_\theta^*H_sA_\theta=\operatorname{diag}\bigl(-(s+j),\,(j-s)\,u_\theta^2\bigr),
 \label{eq:pair-currents-compression}
\end{equation}
because $\langle b,H_sb\rangle=(j-s)u_\theta^2$. The two excited entries differ by $(j-s)v_\theta^2$. For $0<\theta<\pi/4$ and $s\ne j$ two things therefore fail. The prepared plane is not invariant, by Eq.~\eqref{eq:pair-invariance-amplitude}, so $A_\theta$ intertwines no input generator with $H_s$. And the specified input generator differs from the compression.

\subsection{A general mixer and an internal symmetry}
\label{subsec:pair-currents-mixers}

A mixer need not act identically on the two labels. For a complex $2\times2$ matrix $K$ put
\begin{equation}
 T_K=\sum_{i,k}c_i^*K_{ik}d_k,\qquad Y_K=\tfrac1{2i}(T_K-T_K^*),
 \label{eq:pair-currents-mixer}
\end{equation}
so that $Y_I=Y$. Here $T_K$ annihilates $C$ and $T_K^*$ annihilates $D_p$. Expanding the two remaining terms in the four orthonormal states $c_i^*d_k^*\ket{\mathrm{vac}}$ gives
\begin{equation}
 2i\sqrt2\;Y_Kg=\Bigl[(K_{11}-\overline{K_{22}})\,c_1^*d_2^*+(\overline{K_{11}}-K_{22})\,c_2^*d_1^*
 -(K_{12}+\overline{K_{21}})\,c_1^*d_1^*+(K_{21}+\overline{K_{12}})\,c_2^*d_2^*\Bigr]\ket{\mathrm{vac}},
 \label{eq:pair-currents-mixer-image}
\end{equation}
and therefore
\begin{equation}
 \|Y_Kg\|^2=\tfrac14\Bigl(|K_{11}-\overline{K_{22}}|^2+|K_{12}+\overline{K_{21}}|^2\Bigr).
 \label{eq:pair-currents-leakage}
\end{equation}
The vector $Y_Kg$ is orthogonal to $g$. The mixer keeps the ground on its ray exactly when $Y_Kg=0$, that is, when
\begin{equation}
 K=\begin{pmatrix}z&w\\-\overline w&\overline z\end{pmatrix},\qquad z,w\in\C .
 \label{eq:pair-currents-quaternion}
\end{equation}
A nonzero matrix of this form is $r$ times an element of $SU(2)$, with $r^2=|z|^2+|w|^2$. The angle $\theta$ of the prepared state does not enter this criterion, and the criterion does not say whether a laboratory can implement the mixer.

Now suppose that a group acts on the modes through two-dimensional unitary representations, $c\mapsto R_c(h)c$ and $d\mapsto R_d(h)d$ on the columns $c=(c_1,c_2)^T$ and $d=(d_1,d_2)^T$ of annihilation operators. Then $T_K$ goes to $T_{K'}$ with $K'=R_c(h)^*KR_d(h)$. The operators $c_i^*d_k$ and their adjoints are linearly independent, so $Y_K$ is invariant exactly when $R_c(h)K=KR_d(h)$ for every $h$.

\begin{proposition}[Invariant mixers that keep the ground]
\label{prop:pair-currents-invariant-mixer}
\begin{enumerate}
\item[(a)] A nonzero invariant mixer $Y_K$ with $Y_Kg=0$ exists exactly when $R_c$ and $R_d$ are unitarily equivalent.
\item[(b)] The aligned mixer $Y=Y_I$ is invariant exactly when $R_c(h)=R_d(h)$ for every $h$.
\item[(c)] The pair transfer $P_c^*P_d$ is invariant exactly when $\det R_c(h)=\det R_d(h)$ for every $h$.
\end{enumerate}
\end{proposition}
\begin{proof}
For (a), Eq.~\eqref{eq:pair-currents-quaternion} gives $K=rU$ with $r>0$ and $U\in SU(2)$, and invariance makes $U$ a unitary intertwiner, $R_c(h)U=UR_d(h)$. Conversely, if a unitary $V$ satisfies $R_c(h)V=VR_d(h)$, then $V/\sqrt{\det V}$ intertwines, lies in $SU(2)$ and so has the form of Eq.~\eqref{eq:pair-currents-quaternion}. Part (b) is the case $K=I$. For (c), substituting $c\mapsto R_c(h)c$ in $P_c=c_2c_1$ gives $\det R_c(h)\,P_c$, and likewise for $P_d$, so $P_c^*P_d$ is multiplied by $\overline{\det R_c(h)}\det R_d(h)$.
\end{proof}

Condition (c) is weaker than (a). Let the circle group act with charges $(+1,-1)$ on $c_1,c_2$ and $(+2,-2)$ on $d_1,d_2$. Both determinants are one and the pair transfer is invariant. The relation $R_cK=KR_d$ forces $K=0$, because no charge of $c$ equals a charge of $d$. This symmetry allows the pair Hamiltonian and forbids every one-particle mixer between the two orbitals. A charged mediator or reference system could give an invariant joint interaction. It would then belong to the preparation and has to be kept in the model of the measurement. Even an invariant $K$ shows only that the symmetry permits the mixer. Whether a source supplies its Hermitian quadrature $Y_K$ is a further question.

The form of Eq.~\eqref{eq:pair-currents-quaternion} has a simple origin. Over the four modes $a=(c_1,c_2,d_1,d_2)$ the ground is
\[
 g=\frac1{2\sqrt2}\sum_{p,q}\Omega_{pq}\,a_p^*a_q^*\ket{\mathrm{vac}},\qquad
 \Omega=\begin{pmatrix}\epsilon&0\\0&\epsilon\end{pmatrix},\quad
 \epsilon=\begin{pmatrix}0&1\\-1&0\end{pmatrix}.
\]
A one-body unitary acts on $\Omega$ by a congruence $\Omega\mapsto\mathsf R\,\Omega\,\mathsf R^T$ with a unitary $4\times4$ matrix $\mathsf R$. It keeps the ray of $g$ exactly when $\mathsf R\,\Omega\,\mathsf R^T$ is a multiple of $\Omega$, that is, when $\mathsf R$ is a phase times an element of the unitary symplectic group of $\Omega$. The Hermitian generators are the multiples of the identity together with a real Lie algebra of dimension ten: three internal rotations on each orbital and the four-parameter family of Eq.~\eqref{eq:pair-currents-quaternion}. Among the common-label drives, $Y$ belongs to it and $X$, $Z$ do not.

Keeping the ground does not keep the space on which the zero outcome identifies $d$. Take $K=\operatorname{diag}(i,-i)$, which has the form of Eq.~\eqref{eq:pair-currents-quaternion}. Its mixer is the current $X$ taken with opposite signs on the two labels, $Y_K=\tfrac12(c_1^*d_1+d_1^*c_1)-\tfrac12(c_2^*d_2+d_2^*c_2)$. With the vector $w_0\in W$ of Section~\ref{sec:preparation-evolution},
\begin{equation}
 Y_Kg=0,\qquad Y_Ke=-w_0,\qquad Y_Kw_0=-e,\qquad e^{-itY_K}e=\cos t\;e+i\sin t\;w_0 .
 \label{eq:pair-currents-leaving}
\end{equation}
The pulse with $t=\pi/2$ fixes the ground and takes $e$ entirely into $W$. The imbalance then reads zero with certainty although the state has no component on $d$. A mixer admitted by Eq.~\eqref{eq:pair-currents-quaternion} can thus invalidate the reading of Eq.~\eqref{eq:pair-currents-zero}, and the bound on the weight in $W$ used with Eq.~\eqref{eq:pair-ambient-zero} has to be established separately for it. Among the matrices of Eq.~\eqref{eq:pair-currents-quaternion}, only the real multiples of the identity give mixers that map $\Span\{g,e,m\}$ into itself.

\subsection{What the readout and the Hamiltonian generate}
\label{subsec:pair-currents-algebras}

On the input, the generator is $\operatorname{diag}(-(s+j),\,j-s)$ and the input imbalance $A_\theta^*Z_\theta A_\theta$ exchanges $g_\Psi$ and $e_\Psi$. Together they generate all $2\times2$ matrices. On $\Span\{g,e,m\}$, for $0<\theta<\pi/4$ and $s\ne j$, the Hamiltonian $H_s$ has the three distinct eigenvalues $-(s+j)$, $j-s$ and $0$, and $Z_\theta g=u_\theta e+v_\theta m$ has both components nonzero. The spectral projections and these two matrix elements give every matrix unit, so $H_s$ and $Z_\theta$ generate all $3\times3$ matrices. Each ground is cyclic for its algebra, as R2 asks. So is every other nonzero vector. Cyclicity does not select the grounds. The Hamiltonians do. The output space matters here. The three currents and $H_s$ map $\Span\{g,e,m\}$ into itself, and from $g$ they generate exactly this space. If the output is enlarged to include part of $W$, the ground is no longer cyclic for the algebra that $H_s$ and the currents generate.

The ground correlation of the readout shows the two excited lines. Let $Z_\theta(t)=e^{itH_s/\hbar}Z_\theta e^{-itH_s/\hbar}$ be the readout after a wait $t$. Then
\begin{equation}
 \langle g,Z_\theta(t)Z_\theta g\rangle
 =u_\theta^2\,e^{-2ijt/\hbar}+v_\theta^2\,e^{-i(s+j)t/\hbar},
 \label{eq:pair-currents-correlation}
\end{equation}
because $Z_\theta g=b$ and the energies of $e$ and $m$ above the ground are $2j$ and $s+j$. As a function of $t$ this is the transform of a measure with two atoms, and the two dimensions $\Span\{e,m\}$ that the excitation $b$ reaches reproduce it. A realization that also contains the stationary ground whose correlation is measured needs a third dimension: two orthogonal eigenvectors for the two lines, both orthogonal to the ground. The preparation $A_\theta$ prepares $g$ as well as $b$, and its minimal output $\Span\{g,e,m\}$ attains this bound. At $\theta=0$ the weight $v_\theta^2$ vanishes, and at $\theta=\pi/4$ the weight $u_\theta^2$ does. The prepared plane is then $\Span\{g,e\}$ or $\Span\{g,m\}$. It reduces $H_s$, the preparation maps onto this two-dimensional minimal output, and there is no defect. At $s=j$ the two frequencies coincide for every angle.

\subsection{The same currents in the trap}
\label{subsec:pair-currents-orbital}

In the trap of Section~\ref{subsec:orbital-apparatus} the vector potential is that of the symmetric gauge, $\mathbf A(x,y)=\tfrac12\mathcal B\,(-y,x)$. For each particle,
\begin{equation}
 \frac{(\mathbf p-q\mathbf A)^2}{2M}+\frac{M\omega_0^2r^2}{2}
 =\frac{\mathbf p^2}{2M}+\frac{M(\omega_0^2+\Omega_{\rm B}^2)\,r^2}{2}-\Omega_{\rm B}\,(xp_y-yp_x),
 \label{eq:pair-currents-minimal-coupling}
\end{equation}
since $\mathbf p\cdot\mathbf A+\mathbf A\cdot\mathbf p=\mathcal B\,(xp_y-yp_x)$ and $q^2\mathbf A^2=M^2\Omega_{\rm B}^2r^2$. This is Eq.~\eqref{eq:compensated-orbital-parent} with $\omega^2=\omega_0^2+\Omega_{\rm B}^2$. Only the orbital coupling is included. A spin magnetic moment with $g$-factor $g_s$ would add $-g_s\Omega_{\rm B}S_z$. This term commutes with the Hamiltonian and leaves the three singlets where they are. It shifts a level with spin projection $M_S$ by $-g_s\hbar\Omega_{\rm B}M_S$, so the comparison of the ground with the levels of nonzero spin has to be made again. For fixed $a$, $\omega$, $\chi$ and $g_s$, the energy of the lowest component of the shell triplet minus that of $G_\lambda$ is then
\begin{equation}
 \lambda d_a\Bigl(2-\frac{|g_s\chi|}{2}\Bigr)+O(\lambda^2).
 \label{eq:pair-currents-spin-gap}
\end{equation}
For small $\lambda$ the triplet component lies above $G_\lambda$ when $|g_s\chi|<4$ and below it when $|g_s\chi|>4$, and then $G_\lambda$ is not the ground. At $|g_s\chi|=4$ the first order does not decide. The higher levels with spin need the same test.

On the bare shell the currents are familiar observables, and the two kinds of spin separate. The vectors $g,e,m$ have orbital spin one and are singlets of the physical spin. The space $W$ is the physical triplet, and the currents annihilate it. Let $c_i$ and $d_i$ be the valence modes of $p_x$ and $p_y$ with spin $i$. The one-particle angular momentum satisfies $Lp_x=i\hbar p_y$, $Lp_y=-i\hbar p_x$ and $Ls=0$, so on the shell $L=i\hbar\sum_i(d_i^*c_i-c_i^*d_i)=2\hbar Y$. By Eq.~\eqref{eq:pair-currents-action},
\begin{equation}
 Lg=0,\qquad Le=2i\hbar\,m,\qquad Lm=-2i\hbar\,e,
 \label{eq:pair-currents-angular-momentum}
\end{equation}
and the laboratory rotation $R_\varphi=e^{-i\varphi L/\hbar}$ acts on the shell as the rotation $\kappa(\varphi)$ of the pair model. The quadrupoles are
\begin{equation}
 Q_x=\sum_i\frac{x_i^2-y_i^2}{2\ell^2},\qquad Q_y=\sum_i\frac{x_iy_i}{\ell^2}.
 \label{eq:pair-currents-quadrupoles}
\end{equation}
On the three orbitals their only nonzero one-particle matrix elements are $\langle p_x,Q_xp_x\rangle=\tfrac12=-\langle p_y,Q_xp_y\rangle$ and $\langle p_x,Q_yp_y\rangle=\tfrac12=\langle p_y,Q_yp_x\rangle$. Hence
\begin{equation}
 P_{\rm sh}Q_xP_{\rm sh}=Z,\qquad P_{\rm sh}Q_yP_{\rm sh}=X,\qquad P_{\rm sh}LP_{\rm sh}=2\hbar Y
 \label{eq:pair-currents-shell}
\end{equation}
on the shell. Turning the detector through $\varphi$ replaces $Q_x$ by $\cos(2\varphi)Q_x+\sin(2\varphi)Q_y$, which is Eq.~\eqref{eq:pair-currents-readout} again. These are compressions to the bare shell at zero coupling. The detector $Z_0$ of Section~\ref{subsec:orbital-apparatus} is defined on the exact branches through $P_{\rm low}$ and is normalized by $\|P_{\rm low}Q_xG_\lambda\|$. On $\Span\{g,e,m\}$ the two agree at $\lambda=0$.

The splitting of the two lines has a different origin in the two models. For $s\ne j$ the pair Hamiltonian does not commute with $Y$, by Eq.~\eqref{eq:pair-currents-hamiltonian}. The trap Hamiltonian commutes with $L$ at every field. The field separates the angular momenta $\pm2\hbar$, which time reversal keeps degenerate at zero field. In both models the rotation fixes the ground.

Write $m_\lambda=(F_{+,\lambda}-F_{-,\lambda})/(i\sqrt2)$. Then $G_\lambda,e_\lambda,m_\lambda$ is an orthonormal basis of the exact three-dimensional space selected in Section~\ref{subsec:orbital-apparatus}, and it reduces to $g,e,m$ at zero coupling. The zero-angle preparation is the isometry $A$ on $\C^2$ with orthonormal basis $\ket0,\ket1$ given by
\begin{equation}
 A\ket0=G_\lambda,\qquad A\ket1=e_\lambda,\qquad A^*A=I,\qquad \ker A^*=\C m_\lambda
 \label{eq:pair-currents-orbital-preparation}
\end{equation}
within that space. The two grounds are strongly compatible in the sense of Proposition~\ref{prop:ground-polar-compatibility}, with $A\ket0=G_\lambda$ and $A^*G_\lambda=\ket0$. Write $H_\Phi$ for the trap Hamiltonian on this space. Since $H_\Phi e_\lambda=E_2e_\lambda-2i\hbar\Omega_{\rm B}m_\lambda$,
\begin{equation}
 A^*H_\Phi A=\operatorname{diag}(E_0,E_2),\qquad
 (I-AA^*)H_\Phi A\ket1=-2i\hbar\Omega_{\rm B}\,m_\lambda .
 \label{eq:pair-currents-orbital-compression}
\end{equation}
\emph{The Luminous Ground} \cite{LGOrbital} takes this compression as the input Hamiltonian. That is a choice in addition to the prepared plane and the output dynamics. For $\Omega_{\rm B}\ne0$ the plane is not invariant, so the input Hamiltonian is not a restriction of $H_\Phi$ and $A$ does not intertwine the two evolutions. In the pair model the input generator was a restriction and differed from the compression, Eq.~\eqref{eq:pair-currents-compression}.

With $Z_0=\ket{G_\lambda}\bra{e_\lambda}+\ket{e_\lambda}\bra{G_\lambda}$ the statements of Section~\ref{subsec:pair-currents-algebras} carry over. The compression $\operatorname{diag}(E_0,E_2)$ and $A^*Z_0A$ generate all $2\times2$ matrices. For $\Omega_{\rm B}\ne0$ the Hamiltonian has three distinct eigenvalues on the selected space and $Z_0G_\lambda=e_\lambda$ has equal components on $F_{\pm,\lambda}$, so $H_\Phi$ and $Z_0$ generate all $3\times3$ matrices. The ground correlation is
\begin{equation}
 \langle G_\lambda,Z_0(t)Z_0G_\lambda\rangle=e^{-i(E_2-E_0)t/\hbar}\cos(2\Omega_{\rm B}t),
 \label{eq:pair-currents-orbital-correlation}
\end{equation}
with two lines of weight $1/2$ at the frequencies $(E_2-E_0)/\hbar\mp2\Omega_{\rm B}$. The excitation $e_\lambda=Z_0G_\lambda$ alone reaches the two-dimensional space $\Span\{F_{+,\lambda},F_{-,\lambda}\}$, which reproduces this function of $t$. The preparation $A$ also prepares $G_\lambda$. Its minimal output is therefore the whole three-dimensional space.

Without the filter the prepared excitation is $e_Q$ of Section~\ref{subsec:orbital-apparatus}. Its evolving state stays orthogonal to the stationary ground, so the weight outside the prepared plane $\Span\{G_\lambda,e_Q\}$ is one minus the squared return amplitude,
\begin{equation}
 p_\perp(t)=1-|C_0(t)|^2\cos^2(2\Omega_{\rm B}t).
 \label{eq:pair-currents-unfiltered}
\end{equation}
This weight can be spread over several radial levels. Complete departure at $t=\pi/(4|\Omega_{\rm B}|)$ is then not a transfer to one specified line, and it does not ensure a later complete return. Finally, conservation of $L$ does not prevent the departure. The prepared excitation $e_\lambda$ combines two eigenvectors of different energies, and its line is not invariant under their relative phase.
